\section{Spin and orbital angular-momentum flux densities}

For the numerical calculation on the detector, we use angular-momentum flux
densities per unit screen area and per unit positive angular frequency.  The
spin--orbital flux separation follows
Refs.~\cite{barnett_optical_2002,bliokh_conservation_2014}.  As established in
Eq.~\eqref{eq:spectral-intrinsic-angular-momentum-flux}, the exact intrinsic
flux density of $J_3$ in the $x_3$ direction is
\begin{importantbox}
  \begin{align}
    \frac{d\mathcal M^{\mathrm{int}}_{33,+}}{d\omega}
    =(x_1-X_c)\frac{d\mathcal T_{23,+}}{d\omega}
    -(x_2-Y_c)\frac{d\mathcal T_{13,+}}{d\omega},
    \qquad \omega>0.
    \label{eq:one-sided-intrinsic-angular-momentum-flux-density}
  \end{align}
\end{importantbox}

The exact spin flux density is
\begin{importantbox}
\begin{align}
  \frac{d\mathcal M^{\mathrm{spin}}_{33,+}}{d\omega}
  =\frac{4\pi\epsilon_0c^2}{\omega}\operatorname{Im}
  \left[
    {\bf B}(\omega,{\bf x})^*\cdot{\bf E}(\omega,{\bf x})
    -2B_3(\omega,{\bf x})^*E_3(\omega,{\bf x})
  \right],
  \qquad \omega>0.
  \label{eq:one-sided-spin-angular-momentum-flux-density}
\end{align}
\end{importantbox}
The factor $4\pi$ follows from the Fourier normalization in
Eq.~\eqref{eq:Faraday-Fourier-convention-observables} and from using a
one-sided spectrum.  Substitution of
Eq.~\eqref{eq:Fourier-fields-from-Faraday} gives a formula directly in terms
of the available Faraday-tensor components:
\begin{highlightbox}[title={Spin angular-momentum flux density from the Faraday tensor}]
  \begin{align}
    \frac{d\mathcal M^{\mathrm{spin}}_{33,+}}{d\omega}
    =\frac{4\pi\epsilon_0c^3}{\omega}\operatorname{Im}\left[
      -(F^{23})^*F^{10}+(F^{13})^*F^{20}+(F^{12})^*F^{30}
    \right],
    \qquad \omega>0.
    \label{eq:one-sided-spin-angular-momentum-flux-density-Faraday}
  \end{align}
  All components are evaluated at $(\omega,{\bf x})$ on the detector screen.
\end{highlightbox}

The spin angular-momentum flux spectrum is the integral over the detector,
\begin{importantbox}
\begin{align}
  \frac{d\mathcal J^{\mathrm{spin}}_3}{d\omega}
  =\int_A dA\,
  \frac{d\mathcal M^{\mathrm{spin}}_{33,+}}{d\omega}.
  \label{eq:one-sided-spin-angular-momentum-flux}
\end{align}
\end{importantbox}

The intrinsic orbital flux density is defined without introducing a separate
local decomposition of the stress tensor: it is the residual after the spin
part is subtracted from the intrinsic total flux density,
\begin{importantbox}
\begin{align}
  \frac{d\mathcal M^{\mathrm{orb,int}}_{33,+}}{d\omega}
  =\frac{d\mathcal M^{\mathrm{int}}_{33,+}}{d\omega}
  -\frac{d\mathcal M^{\mathrm{spin}}_{33,+}}{d\omega}.
  \label{eq:one-sided-intrinsic-orbital-angular-momentum-flux-density}
\end{align}
\end{importantbox}
Here $d\mathcal M^{\mathrm{int}}_{33,+}/d\omega$ is evaluated from
Eq.~\eqref{eq:one-sided-intrinsic-angular-momentum-flux-density}, with the
stress-tensor components given explicitly in terms of $F^{\alpha\beta}$ by
Eq.~\eqref{eq:one-sided-spectral-linear-momentum-flux-F}.

\begin{highlightbox}[title={Why no simple azimuthal-derivative flux formula is used}]
  For a general field, the exact canonical separation of angular-momentum
  flux contains mixed electric--magnetic terms and a spin--orbit correction.
  Therefore the density expression involving $\partial/\partial\phi$ cannot
  simply be multiplied by $c$ to obtain the flux.  The residual definition
  above is used here because it requires no assumption about the propagation
  direction.
\end{highlightbox}

Its integral gives
\begin{importantbox}
\begin{align}
  \frac{d\mathcal J^{\mathrm{orb,int}}_3}{d\omega}
  =\int_A dA\,
  \frac{d\mathcal M^{\mathrm{orb,int}}_{33,+}}{d\omega}
  =\frac{d\mathcal J^{\mathrm{int}}_3}{d\omega}
  -\frac{d\mathcal J^{\mathrm{spin}}_3}{d\omega}.
  \label{eq:one-sided-intrinsic-orbital-angular-momentum-flux}
\end{align}
\end{importantbox}

In the numerical cases considered here, the longitudinal spectral Poynting
flux is symmetric about the $Ox_3$ axis.  For a detector aperture that is
itself centered on this axis and respects the same symmetry,
\begin{align}
  \frac{d\mathcal S^3_+}{d\omega}(\omega,x_1,x_2)
  &=\frac{d\mathcal S^3_+}{d\omega}(\omega,\rho),
  \qquad \rho=\sqrt{x_1^2+x_2^2},\nonumber\\
  X_c(\omega)&=Y_c(\omega)=0.
  \label{eq:vanishing-energy-flux-centroid}
\end{align}
Equation~\eqref{eq:spectral-extrinsic-angular-momentum-flux} then gives
$d\mathcal J^{\mathrm{ext}}_3/d\omega=0$, irrespective of the values of the
transverse momentum fluxes, and hence
\begin{importantbox}
\begin{align}
  \frac{d\mathcal J^{\mathrm{tot}}_3}{d\omega}
  =\frac{d\mathcal J^{\mathrm{spin}}_3}{d\omega}
  +\frac{d\mathcal J^{\mathrm{orb,int}}_3}{d\omega}.
  \label{eq:angular-momentum-split-zero-extrinsic}
\end{align}
\end{importantbox}
Thus the vanishing extrinsic contribution follows from the symmetry of the
measured longitudinal energy-flux distribution, not from an assumed
vanishing of the transverse momentum fluxes.  Equation~\eqref{eq:one-sided-intrinsic-orbital-angular-momentum-flux-density}
remains the definition used for spatial flux-density maps.
